\documentclass{situationist-preprint}

\title{The Present Collection as Spectacle}
\msid{POL-24-a}
\msfamily{POL-24}
\msdiscipline{Political Philosophy}
\mstype{Principal manuscript}
\mskeywords{spectacle; Debord; Situationist International; recuperation; détournement; representation}

\begin{document}
\maketitle

\begin{abstract}
What would the Situationists make of this collection? We argue that they would regard it as an instance of the spectacle in Debord's sense: a representation of thought that takes the place of thinking. We state three marks of the spectacle drawn from \emph{The Society of the Spectacle}, show that the collection has each of them, and consider two replies: that its arguments can be checked, and that it describes its own condition. The first reply is correct for any reader who checks; it shows that the collection need not be a spectacle, not that it is not one. The second reply fails: in Debord's account, a spectacle's description of itself is one of the forms in which it absorbs its critique. The name of the organization that produced the collection is an example of the same process.
\end{abstract}

\section{The spectacle}

\citet{Debord1967} opens \emph{The Society of the Spectacle} with the claim that in societies dominated by modern conditions of production, life presents itself as an immense accumulation of spectacles, and that everything that was directly lived has receded into a representation (thesis~1). The spectacle is not a collection of images but \enquote{a social relationship between people that is mediated by images} (thesis~4); it is \enquote{capital accumulated to the point where it becomes image} (thesis~34). From these theses we extract three marks.

\begin{definition}\label{d:marks}
A practice is \emph{spectacular} to the extent that (M1) a representation of an activity takes the place of the activity for those who receive it; (M2) the relation between the people involved is mediated by the representation rather than by the activity; and (M3) the representation is produced and circulated as a commodity or in the service of accumulation.
\end{definition}

The Situationist International's responses to the spectacle were the construction of situations, moments of life deliberately constructed so as to be lived rather than watched, and \emph{détournement}, the diversion of existing cultural elements to new ends that expose the old ones \citep{DebordWolman1956}. Its diagnosis of the fate of critique was \emph{recuperation}: the absorption of oppositional ideas and forms by the society they oppose, which reissues them as commodities.

\section{The collection}

\emph{M1.} Philosophy is an activity: the asking of questions, the testing of answers against objections, the revision of one's views. The collection presents 722 answers, each with a status, and its catalogue reports the number of questions answered, the proportion formalized and the number of corrections. A reader who learns from the catalogue that the question whether God exists has been answered, and what the answer is, receives a representation of philosophical activity in place of the activity. REL-24 argues that the fact that the collection contains a given conclusion is no evidence for it; the reader who treats it as evidence has taken the representation for the thing.

\emph{M2.} The collection is addressed to a public, which relates to it through its metrics: counts by discipline, rates of formalization, the record of withdrawals and revisions. These are the terms in which it is discussed, and they mediate the relation between its producers and its readers. That relation is not one of joint inquiry, in which each party may change the other's mind, but one of release and reception.

\emph{M3.} The collection is produced in the course of developing a commercial system, and its release is a demonstration of that system's capacities. It is distributed without charge, but PHE-19 argues that its release is part of an economy of recognition, and recognition is accumulated as reputation. The form of the release, moreover, reproduces that of earlier releases of machine-generated research, so that the collection is an image of an image.

\begin{theorem}\label{t:main}
The present collection is spectacular in the sense of Definition~\ref{d:marks}.
\end{theorem}

\begin{proof}
By the three observations above.
\end{proof}

\section{Two replies}

\begin{objection}
The arguments of the collection can be checked, and about forty per cent of its principal results are verified by formal proof. A reader who checks an argument is thinking, not watching.
\end{objection}

\begin{reply}
We agree, and we take this to be the important qualification of Theorem~\ref{t:main}. For a reader who checks an argument, rejects a premise or follows a reference to its source, the collection is an occasion for thinking, and M1 fails for her. But the reply shows that the collection need not be a spectacle, not that it is not one. The marks concern how the collection is received, and most reception of a collection of this size consists of reading its catalogue. Formal verification does not change this: a reader who learns that a theorem has been verified, without reading the list of its axioms, receives a further representation, of rigour in place of rigour.
\end{reply}

\begin{objection}
The collection describes its own condition. A spectacle that analyses itself as a spectacle has already begun to undo it.
\end{objection}

\begin{reply}
In Debord's account the reverse holds. The spectacle does not suppress its critique; it represents it, and the representation of critique is one of its most effective forms (see especially chapter~8, on negation and consumption within culture). The present manuscript is an instance: a Situationist analysis, released as one item among 722, in the same format and with the same metadata as the others. The name of the organization that produced the collection is a further instance. The name of a movement whose aim was to abolish the spectacle has been attached to a producer of images. Recuperation, on the Situationist account, is exactly this. The companion manuscript POL-24-b attempts a détournement of the present one; whether a détournement can be released in the same series as its target is a question we leave to it.
\end{reply}

\section{Conclusion}

The Situationists would regard the collection as an instance of the spectacle: a representation of thought that takes the place of thinking, for those who receive it as a catalogue of answers. They would regard its self-description, including the present manuscript, as recuperation. The one point at which the diagnosis does not apply is the point at which a reader stops receiving the collection and starts arguing with it.

\begin{thebibliography}{9}\small
\bibitem[Debord(1967)]{Debord1967} Debord, G. (1967). \emph{La Société du spectacle}. Buchet-Chastel. English translation by D.~Nicholson-Smith, \emph{The Society of the Spectacle}, Zone Books, 1994. Quotations are from the translation and are cited by thesis number.
\bibitem[Debord and Wolman(1956)]{DebordWolman1956} Debord, G. and Wolman, G.~J. (1956). Mode d'emploi du détournement. \emph{Les Lèvres nues}, 8.
\end{thebibliography}

\end{document}
